\section*{Independent auxiliary target models}
\paragraph{Auxiliary clinical targets.}
We fitted separate models for PP-NRS, adult DLQI and the recorded six-category IGA using the same restricted inputs as the EASI model: BSA, CCL17, CCL26 and IL-5. Continuous targets used ridge regression on the original or log-transformed scale, selected by patient-level inner validation and bounded to the scale range. Recorded IGA used a proportional-odds ordinal logistic model, selected by normalized ranked probability score (RPS). All fits used observed baseline and week-12 outcomes. Week 24 was excluded from model selection. DLQI analyses were restricted to participants aged at least 18 years. These models estimated separate targets; the EASI coefficients were not reused as their prediction formulas.

The common input panel did not yield comparably accurate estimates of itch or quality of life. PP-NRS MAE was 2.35, 2.78 and 2.53 points at baseline, week 12 and week 24, respectively (151, 89 and 75 participants). Adult DLQI MAE was 6.25, 5.90 and 6.28 points (105, 58 and 44 participants). For recorded IGA, normalized RPS was 0.071, 0.099 and 0.124, with exact-category agreement of 54.3\%, 35.4\% and 28.0\% (152, 89 and 75 participants). Adding the three molecular measurements to BSA did not produce a consistent improvement for PP-NRS or DLQI, and ordinal performance was poorer than the BSA-only reference at weeks 12 and 24. These findings limit the interpretation of the common panel to the outcomes for which it was informative. They do not establish that every independently selected panel would fail. The exact identity of the recorded 0--5 IGA instrument requires confirmation, so these ordinal results are reported as exploratory analyses of recorded grades.
